\section{Fourth cumulants of the limiting components}
\label{sec:ultraviolet-limits}

The fourth-cumulant estimate for a mixed limiting law does not by
itself control each ergodic component. We use the stronger estimate
on products of translated Wick polynomials. It bounds their
translation-invariant conditional expectations, and hence the
second and fourth moments of almost every component.

\subsection{From distinct-point products to the invariant part}

Let $U(y)$,
$y\in\Z^d$, be the operators of a measure-preserving action on
a probability space, where $d\ge1$, with $U(y)G=G\circ T_{-y}$. Let
$P_{\mathrm{lat}}$ be the orthogonal projection in $L^2$ onto
the vectors invariant under this action.

\begin{lemma}\label{lem:UV-invariant-product-bound}
Suppose that $G$ is real and belongs to every $L^p$, $p<\infty$.
If $u\ge0$ and
\begin{equation}\label{eq:UV-distinct-product-bound}
 \left|\E\prod_{a=1}^J U(y_a)G\right|\le u^J
\end{equation}
for every finite family of pairwise distinct points of $\Z^d$,
then $|P_{\mathrm{lat}}G|\le u$ almost surely.
\end{lemma}

\begin{proof}
For $B_n=\{0,\ldots,n-1\}^d$ put
$A_n=n^{-d}\sum_{y\in B_n}U(y)$. These contractions converge
strongly in $L^2$ to $P_{\mathrm{lat}}$.
Indeed, the orthogonal complement of the invariant space is the
closed span of $(U(y)-1)H$. On such a vector, cancellation of
the common terms in $B_n$ and $y+B_n$ gives
\[
 \|A_n(U(y)-1)H\|_2\le2|y|_1n^{-1}\|H\|_2.
\]
On invariant vectors the averages are the identity, proving
the asserted convergence by the orthogonal decomposition.

For an even integer $J\ge2$, expand $\E(A_nG)^J$.
The distinct-index terms obey \eqref{eq:UV-distinct-product-bound}.
There are at most $\binom J2 n^{d(J-1)}$ remaining tuples:
each has a repeated pair of indices. H\"older's inequality
bounds their individual expectations in modulus by
$\|G\|_J^J$. Therefore
\[
 \E|A_nG|^J
       \le u^J+\binom J2 n^{-d}\|G\|_J^J.
\]
Along a subsequence the $L^2$ convergence is almost sure.
Fatou's lemma gives $\E|P_{\mathrm{lat}}G|^J\le u^J$
for every even $J$. For $t>u$, Markov's inequality now gives
$\Pr(|P_{\mathrm{lat}}G|>t)\le(u/t)^J$.
Let $J\to\infty$ and then $t\downarrow u$ through a countable
sequence.
\end{proof}

Let a translation-invariant law $\mu$ on tempered fields have
the decomposition $\mu=\int\mu_\xi\,\dd\nu(\xi)$.
Write $I$ for its translation-invariant sigma algebra and
$P=\E_\mu[\,\cdot\mid I]$. Apply the lemma to translations
on $\lambda\Z^3$, with $\lambda>0$ fixed. The range of $P$
is contained in that of $P_{\mathrm{lat}}$, so
$P=PP_{\mathrm{lat}}$. Positivity of conditional expectation
therefore gives
\begin{equation}\label{eq:UV-component-product-bound}
 |PG|\le P|P_{\mathrm{lat}}G|\le u,\qquad
 |\mu_\xi(G)|\le u\quad\text{for }\nu\text{-almost every }\xi.
\end{equation}
The second assertion uses the conditional-law representation of
the decomposition. Integrability under the conditional laws and
the representation for this integrable $G$ follow by truncating
its positive and negative parts.

\subsection{Passing the Wick-product estimate to a limit law}

Fix the charge $e>0$ and a real compactly supported test $f$.
Use the Wick polynomials $W_m$ in \eqref{eq:vtm-Wick},
with variance argument $v\ge0$ and $m=2,4$.

Take a joint cutoff sequence defining a limit law $\mu$.
Suppose that the lattice estimates give compactly
supported real tests $f_i$, all supported in one fixed cube,
and real numbers
\[
 v_i,\ c_{2,i},\ c_{4,i},\ u_{2,i},\ u_{4,i},
\]
with
$v_i\ge0$, such that
\begin{equation}\label{eq:UV-test-convergence}
 \|f_i-f\|_2\to0,\qquad
 v_i\to\|f\|_2^2,\qquad
 \sup_i(|c_{2,i}|+|c_{4,i}|)<\infty.
\end{equation}
For a fixed $\lambda>0$, assume that
\begin{equation}\label{eq:UV-cutoff-product-bound}
 \left|\E_i\prod_{y\in Y}
       W_m(F^\sharp(\tau_y f_i);v_i,c_{m,i})\right|
                     \le u_{m,i}^{|Y|},\qquad m=2,4,
\end{equation}
for every finite $Y\subset\lambda\Z^3$ and all sufficiently
large $i$ depending on $Y$. This is the translated-cell form
of the transfer estimate; the volume condition makes the
points of $Y$ incongruent modulo the torus.
The reference estimates also give
\begin{equation}\label{eq:UV-reference-bounds}
 \begin{split}
 0\le u_{4,i}&\le Ce^\sigma,\qquad
 0\le(1+|c_{2,i}|)u_{2,i}\le Ce^\sigma,\\
 |c_{4,i}-3c_{2,i}^2-b(f)|
                  &\le C(e^\sigma+\eta_i^{1/4}),
 \end{split}
\end{equation}
with $C$ independent of the cutoff, $\sigma>0$ and $0<e\le1$.
For the chamber laws, Proposition~\ref{prop:transfer-UV-inputs}
supplies these hypotheses throughout the charge interval
\eqref{eq:constants-probe-charge}.

Pass to a subsequence on which the four bounded sequences
$c_{2,i},c_{4,i},u_{2,i},u_{4,i}$ converge, with limits
$c_2,c_4,u_2,u_4$. The weak limit remains $\mu$.
Equation \eqref{eq:UV-reference-bounds} gives
\begin{equation}\label{eq:UV-reference-limit}
 u_4\le Ce^\sigma,\quad
 (1+|c_2|)u_2\le Ce^\sigma,\quad
 |c_4-3c_2^2-b(f)|\le Ce^\sigma.
\end{equation}
Set $G_m=W_m(F(f);\|f\|_2^2,c_m)$.
Then \eqref{eq:UV-cutoff-product-bound} passes to
\begin{equation}\label{eq:UV-law-product-bound}
 \left|\E_\mu\prod_{y\in Y}U(y)G_m\right|\le u_m^{|Y|}
                   \qquad(Y\subset\lambda\Z^3\text{ finite}).
\end{equation}

To justify this passage for the varying tests, expand the finite
product of Wick polynomials into monomials. For a monomial
of positive total degree $J$, telescope the replacement of each
$F^\sharp(\tau_y f_i)$ by $F^\sharp(\tau_y f)$.
All these tests have support in one open $R_i$-cube once $i$
is large. The test-map variance and product-moment bounds
then give an error at most
\[
 Jc_J B^{J-1}\|f_i-f\|_2,
 \qquad B=\sup_i\max(\|f_i\|_2,\|f\|_2).
\]
This tends to zero. The moments of the fixed tests converge
by Section~\ref{sec:continuum-limits}. The polynomial coefficients
converge by the selected coefficient subsequence and the
variance convergence in \eqref{eq:UV-test-convergence}.
There are only finitely many monomials for each fixed $Y$.
This proves \eqref{eq:UV-law-product-bound}.

\subsection{The cumulant in an ergodic component}

\begin{proposition}\label{prop:UV-component-cumulant}
Under \eqref{eq:UV-test-convergence}--\eqref{eq:UV-reference-bounds},
suppose that almost every translation-ergodic component of $\mu$
satisfies Gaussian domination. Then
\begin{equation}\label{eq:UV-component-cumulant}
 \left|\kappa_4^{\mu_\xi}(f)-e^4b(f)\right|
                 \le(8C+3C^2)e^{4+\sigma}
\end{equation}
for almost every component. In particular, if $b(f)>0$ and
$(8C+3C^2)e^\sigma\le b(f)/2$, its fourth cumulant is positive.
\end{proposition}

\begin{proof}
Gaussian domination of $\mu$ makes the real polynomials $G_2,G_4$
belong to every $L^p(\mu)$. Apply
Lemma~\ref{lem:UV-invariant-product-bound} to
\eqref{eq:UV-law-product-bound}, and then use
\eqref{eq:UV-component-product-bound}. On one set of full
component measure,
\[
 |d_m|\le u_m,\qquad d_m=\E_{\mu_\xi}G_m,\qquad m=2,4.
\]
Intersect with the full-measure set of components having Gaussian
domination. For a component in this intersection, put $X=F(f)$
and $v=\|f\|_2^2$. Gaussian domination gives $\E_{\mu_\xi}X=0$
and all moments finite. Thus
\[
 \kappa_4^{\mu_\xi}(f)
  =\E_{\mu_\xi}H_4(X;v)
          -3\bigl(\E_{\mu_\xi}H_2(X;v)\bigr)^2.
\]
The definition of $G_m$ gives the exact identity
\begin{equation}\label{eq:UV-cumulant-identity}
 e^{-4}\kappa_4^{\mu_\xi}(f)
                =c_4+d_4-3(c_2+d_2)^2.
\end{equation}
Consequently \eqref{eq:UV-reference-limit} implies
\begin{align*}
 |e^{-4}\kappa_4^{\mu_\xi}(f)-b(f)|
 &\le |c_4-3c_2^2-b(f)|+u_4+6|c_2|u_2+3u_2^2\\
 &\le8Ce^\sigma+3C^2e^{2\sigma}
 \le(8C+3C^2)e^\sigma.
\end{align*}
This proves the assertion uniformly over the full-measure set
of components. A centered normal variable has zero fourth
cumulant, so positivity also proves that these laws are
non-Gaussian.
\end{proof}




\subsection{Completion for the prescribed volume-first law}

\begin{proof}[Proof of \Cref{thm:vtm-main}]
The scale calculation in \cref{sec:vtm-bounds} gives
\eqref{eq:transfer-W3} and the reference volume error.
The reference and free-coefficient hypotheses are exactly
\eqref{eq:transfer-terminal-input} and \eqref{eq:transfer-free-input}.
Thus \cref{prop:transfer-UV-inputs}, with the uniform constant
\eqref{eq:transfer-CU}, proves
\eqref{eq:vtm-main-errors}--\eqref{eq:vtm-main-finite-cumulant}.

For the prescribed parent, use the finite-torus diagonal
in \cref{sec:continuum-limits}. Set
$f_i=f^{\eta_i}$ and $v_i=v_{\eta_i}$.
The compact support, test convergence and variance convergence
are \eqref{eq:cells-smooth-test},
\eqref{eq:cells-test-error} and \eqref{eq:cells-variance-limit}.
At this fixed charge, $\lambda=R^*/L$ is fixed.
Each finite $Y\subset\lambda\Z^3$ eventually consists of
distinct torus translates. The four reference sequences have
a convergent subsequence, whose weak limit is still the
prescribed $\mu$. The product passage proved above therefore
applies to this same law, with $C=C_U$.
Its components have Gaussian domination by the explicit
hypothesis \eqref{eq:vtm-component-G}.
Apply \cref{prop:UV-component-cumulant} and impose
\eqref{eq:constants-probe-charge} when $b(f)>0$.

If a separate argument gives other properties on a full set
of components of this parent, intersect that set with the
one obtained here. The component laws themselves have not
been replaced.
\end{proof}
